% !TEX root = ../main.tex

\chapter{绪论}

\section{研究背景与意义}

随着具身智能、移动机器人和深度学习技术的快速发展，机器人正在从结构化、封闭式的实验室环境走向消防巡检、室内导览、园区巡逻、地下管廊检查和灾后排查等真实应用场景。在这些场景中，机器人不仅要感知周围环境，还要在地图不完备、光照变化明显、有障碍物、通道空间受限以及定位先验较弱的条件下完成自主移动。导航因此不再只是移动机器人“从一点到另一点”的路径计算问题，而是机器人完成巡检、导览、搬运、搜救和服务任务之前必须具备的基础能力。没有稳定的导航能力，感知结果难以转化为有效行动，任务规划也无法落到具体空间执行上。因此，导航逐渐成为连接视觉感知、任务理解、局部决策和底层执行的核心环节\cite{pei2026embodied_nav}。

从应用需求来看，让机器人进入上述场景具有直接的现实意义。消防巡检、灾后排查和地下管廊检查等任务往往伴随烟尘、高温、有毒气体、坍塌风险或长时间重复劳动，完全依赖人工会带来较高安全成本和人力消耗；室内导览、园区巡逻、楼宇安防和场馆服务等任务虽然危险性较低，但也对连续运行、路径复现和场景适应能力提出了要求。与人相比，机器人可以在危险或重复环境中保持稳定执行，并通过视觉传感器持续记录环境状态，因此适合承担移动感知与目标到达类任务。然而，不同机器人平台对环境的适应能力和工程部署条件存在明显差异，因此有必要进一步明确本文所关注的机器人形态。

从机器人平台类型来看，不同形态的移动机器人具有不同的适用边界。人形机器人在结构上更接近人类，理论上能够适应面向人设计的环境，并完成一定的人机交互任务，但其系统结构复杂、控制难度高，稳定性和工程部署成本仍然较高。轮式底盘机器人结构简单、控制成熟、能耗较低，在平整地面和规则室内环境中具有较好的应用基础，但其运动能力较依赖地面条件，面对台阶、障碍物、碎石、坡道或非结构化地形时通过性受到限制。相比之下，四足机器人在机动性、稳定性和环境适应性之间具有较好的平衡：其足式运动方式能够跨越一定高度的障碍物，适应复杂地面和狭窄空间，同时相比人形机器人具有更低的结构复杂度和更高的工程可部署性。因此，在灾害现场、工业巡检、园区巡逻和室内外混合环境等任务中，四足机器人是一类具有代表性的具身导航平台\cite{yang2019quadruped_review,hwangbo2019anymal,lee2020learning}。

基于上述分析，如何使四足机器人根据简单、直观的目标表达，在真实环境中自主到达指定位置，是具身导航研究面向工程落地时必须回答的问题。传统移动机器人导航通常依赖全局坐标、度量地图和预设航点。这一路线在仓储物流、工厂搬运、医院配送、商场清洁和室内导览等平整、规则环境中已经形成较成熟的系统链路，轮式平台也在这些场景中承担了重复运输、固定路线巡检和服务引导等角色。然而，在临时巡检、陌生室内空间或弱 GPS 场景中，精确地图往往难以及时获取，普通用户也不容易用坐标描述目标位置。在真实应用中，操作者往往更容易通过目标图像、路线图像或关键场景记忆表达“想让机器人到达哪里”。

视觉目标导航为这一问题提供了更自然的任务表达方式：操作者只需提供目标位置的第一视角图像，或由若干关键图像组成的拓扑地图，机器人便可根据当前图像与目标图像之间的关系生成局部运动。这种“看见某处并走过去”的形式更贴近人的空间记忆方式，也更适合导览复现、巡检路线记录和弱先验环境中的目标到达任务\cite{zhu2017targetdriven,shah2021ving}。但是，对于四足平台而言，并不是将轮式导航算法直接更换载体即可使用。四足机器人的步态运动会引入相机抖动和机体俯仰，速度指令与真实位移之间也会存在非线性偏差，高层视觉策略还必须与底层运动主机、通信协议和安全停机机制形成稳定闭环。

近年来，通用视觉导航模型和生成式策略的发展为上述问题提供了新的技术基础。通用导航模型（General Navigation Model, GNM）和视觉导航 Transformer（Visual Navigation Transformer, ViNT）将历史图像观察、目标条件和局部航点预测纳入统一模型\cite{gnm2022,vint2023}；目标掩码扩散导航模型（Navigation with Goal Masked Diffusion, NoMaD）进一步引入动作扩散策略，通过条件去噪采样生成多条局部轨迹候选，并利用目标掩码机制（goal mask）统一目标导航与探索行为\cite{nomad2023,chi2025diffusion}。与直接回归单一路径相比，扩散式轨迹生成能够描述多模态决策分布，更适合走廊分叉、障碍绕行和目标不确定等情形。然而，现有公开验证仍多集中在轮式平台或离线数据集上，面向四足机器人真实部署时的实时推理、视觉扰动、运动接口适配和安全闭环问题仍需要系统研究。

基于上述背景，本文以 NoMaD 为高层视觉导航模型，以云深处 Lite3 四足机器人为部署对象，研究在室内狭窄、弱先验、无全局精确地图的环境中，如何将扩散式视觉目标导航模型部署到资源受限且运动扰动明显的四足平台上，使其能够根据真实相机图像与目标图像形成可实时运行、可仿真验证、可真机执行的闭环导航系统。

\begin{figure}[htbp]
    \centering
    \includegraphics[width=0.5\linewidth]{intro-background-placeholder.png}
    \caption{四足机器人视觉目标导航目标场景示意图}
    \label{fig:intro-background-placeholder}
\end{figure}

围绕这一问题，本文的研究意义主要体现在三个层面。首先，在应用层面，本文将视觉目标导航与四足平台结合，使导航任务不再停留于离线图像匹配或仿真指标，而是面向室内巡检、导览和受限空间通行等具体场景。其次，在算法层面，本文关注动作扩散策略在边缘设备上的推理效率和轨迹质量，通过去噪扩散隐式模型（Denoising Diffusion Implicit Model, DDIM）采样加速、无分类器引导（Classifier-Free Guidance, CFG）、测试时扩展（Test-Time Scaling, TTS）候选筛选和视觉编码器替换等实验，分析扩散式导航模型在真实部署中的速度、质量和稳定性权衡。最后，在系统层面，本文构建从统一推理模块、状态机、导航主机、中层控制映射到 Lite3 UDP 运动接口的分层闭环架构，使高层视觉导航策略能够接入 MuJoCo 仿真环境和 Jetson Orin 真机平台，为后续四足视觉导航研究提供可复现的工程基础。

\section{国内外研究现状}

四足机器人视觉目标导航是一个交叉性较强的研究问题，涉及传统导航与视觉同步定位与建图（Simultaneous Localization and Mapping, SLAM）、学习型视觉导航、拓扑记忆、生成式动作策略以及四足机器人真实部署等多个方向。为了厘清相关技术的发展脉络并明确本文的研究切入点，本节按照“导航建模—目标表达—扩散策略—四足部署”的逻辑，对国内外研究现状进行综述。

\subsection{传统导航、视觉 SLAM 与学习型导航}

面向已知或可建图环境中的机器人自主移动问题，传统导航方法通常由状态估计、SLAM、路径规划和运动控制等模块组成。Thrun 等\cite{thrun2005probabilistic}围绕不确定感知与运动问题，系统化建立了概率机器人学框架，使定位、建图和决策能够在统一的概率模型下描述。该框架具有清晰的理论结构和可解释性，但在真实场景中往往依赖较准确的传感器模型与环境假设。Hess 等\cite{cartographer2016}针对二维激光实时建图提出 Cartographer，通过子图匹配和回环检测提升了工程可用性，在激光雷达配置充足的移动平台上表现稳定；当机器人仅依赖单目视觉或处于纹理变化明显、遮挡频繁的室内环境时，该类方法的直接适用性会受到限制。Campos 等\cite{orbslam3_2021}提出的定向快速特征与旋转二进制鲁棒特征同步定位与建图系统（Oriented FAST and Rotated BRIEF-Simultaneous Localization and Mapping 3, ORB-SLAM3）面向视觉、视觉惯性和多地图 SLAM 问题，在特征匹配、局部建图和回环优化方面形成了成熟框架，能够获得较高定位精度；但其性能仍与视觉特征质量、标定精度和几何一致性密切相关。

国内学者也对定位建图与导航技术进行了持续研究。毛军等\cite{mao2022slam_review}围绕惯性、视觉和激光雷达 SLAM 的技术演进，梳理了多传感器融合方法的适用条件和工程链路，为理解不同传感器组合下的定位建图方案提供了较完整的综述。该类研究对导航系统的底层定位能力具有重要参考价值，但关注重心仍在环境表示与位姿估计，而不是由目标图像直接驱动的端到端导航决策。裴凌和熊超然\cite{pei2026embodied_nav}从具身导航的概念、方法和挑战出发，强调导航研究正在从几何到达转向任务驱动、语义理解和自主决策，这一判断与本文关注的弱先验视觉目标导航具有一致性；不过，这篇文章并未对具体四足平台给出完整部署链路。

随着深度学习的发展，研究者开始尝试绕开完整几何建模，直接从视觉观察中学习导航策略。Zhu 等\cite{zhu2017targetdriven}针对室内目标驱动视觉导航问题，利用深度强化学习建立当前图像、目标图像与动作之间的映射，表明图像目标任务可以通过学习方法求解。该工作打开了视觉目标导航的研究空间，但仿真依赖较强，距离真实机器人稳定部署仍有差距。Anderson 等\cite{anderson2018evaluation}面向具身导航评价不统一的问题提出标准化评价指标，推动了导航任务评测方式的规范化；然而，评价指标解决的是“如何比较”，并不能替代真实平台上的感知、控制和安全闭环设计。总体来看，传统方法在结构化环境中可靠，学习型方法在任务表达和跨场景泛化方面更灵活，面向四足机器人弱先验部署时仍需要将两者的工程稳定性与视觉决策能力结合起来。

\subsection{视觉目标导航与拓扑记忆}

视觉目标导航关注机器人如何根据目标图像、目标对象或其他视觉线索生成移动行为。相较于点目标导航，图像目标任务减少了对全局坐标的依赖，更接近普通操作者在真实场景中的目标描述方式。Wijmans 等\cite{ddppo2020}针对点目标导航的大规模训练问题，基于分布式近端策略优化（Proximal Policy Optimization, PPO）构建了高效训练框架，并在仿真环境中取得较高成功率。这一工作说明大规模交互训练可以显著提升导航策略性能，但坐标目标与本文关注的目标图像/拓扑地图节点仍存在差异。陈铂垒等\cite{chen2025object_nav_review}围绕物体目标导航系统总结了语义地图、强化学习和大模型辅助导航等方向，为理解物体目标任务提供了较新的国内综述视角；物体类别目标更强调“找到某类物体”，而本文的目标图像导航更强调“到达某个视觉位置”。司马双霖等\cite{sima2023vln_review}系统梳理了视觉语言导航中语言指令、视觉感知和路径规划的结合方式，说明多模态目标表达具有广阔应用潜力；但语言指令天然存在歧义，通常还需要复杂的语义解析模块，部署难度高于目标图像驱动方式。

为了降低对全局精确地图的依赖，拓扑记忆方法将长程导航拆分为一系列局部视觉子目标。Savinov 等\cite{savinov2018sptm}提出半参数拓扑记忆（Semi-Parametric Topological Memory, SPTM），通过视觉相似度和局部策略沿记忆图前进，避免了完整几何地图重建。该方法适合用稀疏视觉节点组织长程路线，但对记忆节点质量和局部控制策略仍较敏感。Chaplot 等\cite{chaplot2020nts}将学习表示与拓扑规划结合，提出神经拓扑 SLAM，在复杂室内场景中同时构建拓扑图并规划路径；这种框架增强了规划能力，也带来了较高系统复杂度。Shah 等\cite{shah2021ving}面向开放世界视觉目标导航提出开放世界视觉目标导航方法（Learning Open-World Navigation with Visual Goals, ViNG），通过视觉路标和目标图像实现真实环境长程导航，展示了视觉记忆路线在实际机器人中的价值；其公开系统主要基于轮式平台，面对四足步态扰动、速度映射和安全停机等问题时仍需要额外设计。

本文采用的拓扑地图思路与上述拓扑记忆方法一脉相承，但研究重点并非重新提出拓扑导航算法，而是将有序视觉节点作为四足机器人真实部署中的任务载体。系统在滑动窗口中定位当前节点，选择实时子目标图像，再由动作扩散策略生成局部航点。这样既保留了视觉目标导航的直观性，又降低了对全局地图和精确定位的依赖，使任务表达能够与四足机器人在真实室内环境中的执行链路衔接起来。

\subsection{动作扩散策略与通用视觉导航模型}

扩散模型最初在图像生成领域取得突破，其核心思想是通过逐步加噪和反向去噪学习复杂数据分布。Ho 等\cite{ddpm2020}提出去噪扩散概率模型（Denoising Diffusion Probabilistic Model, DDPM），为高质量生成建模提供了重要基础，但逐步采样也带来了较高计算开销。Song 等\cite{ddim2020}提出 DDIM，通过非马尔可夫反向过程减少采样步数，使扩散模型在不重新训练的情况下获得更快推理速度；采样步数压缩过度时，生成质量和稳定性仍可能下降。Chi 等\cite{chi2025diffusion}将扩散思想系统引入机器人视觉运动控制，提出 Diffusion Policy，将动作序列建模为条件生成过程，并结合滚动时域控制实现闭环执行。该方法能够表达多模态操作轨迹，但在实时控制中必须处理迭代采样、部署延迟和边缘设备算力约束等问题。

在视觉导航领域，通用导航模型逐渐形成了从轨迹回归到生成式预测的发展脉络。Shah 等\cite{gnm2022}提出 GNM，尝试用统一模型处理不同机器人平台上的图像观察和目标条件，说明跨平台导航策略具有可行性。ViNT\cite{vint2023}进一步利用 Transformer 融合历史观察和目标图像，提升了时序视觉表征能力。Sridhar 等\cite{nomad2023}提出 NoMaD，将 goal mask 与动作扩散策略结合，在同一模型中统一目标导航和无目标探索，并生成多条候选轨迹。这类模型为视觉目标导航提供了更强的局部决策表达能力，尤其适合多路径选择和局部目标不确定的场景；但公开实验仍主要围绕轮式平台展开，扩散采样在 Jetson Orin 等边缘设备上的实时性，以及生成航点到四足底层速度指令的映射问题仍有待解决。

因此，动作扩散策略为四足视觉目标导航提供了有吸引力的高层模型，但它并不自动完成真实部署。对于 Lite3 这类四足平台，系统还必须关注采样步数如何压缩、候选轨迹如何筛选、视觉编码器替换是否影响导航质量、速度指令如何下发到底层运动主机、异常输出如何触发安全停机等问题。

\subsection{四足机器人视觉导航与真实部署}

四足机器人研究长期围绕复杂地形通过、动态平衡和感知运动控制展开。杨钧杰等\cite{yang2019quadruped_review}综述了四足机器人机构、步态规划、控制方法和应用方向，较早整理了国内外四足机器人研究脉络；由于该方向当时仍主要聚焦运动控制，对视觉目标导航和生成式策略的讨论相对有限。Hwangbo 等\cite{hwangbo2019anymal}基于强化学习训练 ANYmal 完成敏捷运动，展示了学习控制在真实四足平台上的潜力。Lee 等\cite{lee2020learning}进一步面向复杂地形行走提出学习型运动控制方法，提升了四足机器人在野外地形中的鲁棒性。这些工作强化了四足平台的底层运动能力，但导航目标理解、视觉子目标选择和高层决策仍需要额外模块配合。

在感知导航方面，Miki 等\cite{miki2022perceptive}通过外感知与本体感知融合，使四足机器人能够在自然地形中实现鲁棒穿越，体现了感知信息对足式运动的重要价值。Ren 等\cite{ren2023hvn}面向四足机器人导航中的落足点适应问题，提出分层视觉导航与落足点适应学习相结合的系统，由国内团队完成了视觉导航与四足落足控制的结合。该工作与本文同样关注四足导航部署，但其目标表达和高层策略形式不同，本文更关注目标图像条件下的扩散式局部轨迹生成。Kareer 等\cite{kareer2023vinl}提出视觉导航与运动控制系统（Visual Navigation and Locomotion over Obstacles, ViNL），将视觉导航策略与视觉运动控制策略协同部署到室内四足平台上，使机器人能够在跨越障碍的同时完成导航任务；该系统验证了四足视觉导航的可行性，但依赖特定训练和感知设置，并未展开 NoMaD 类通用扩散导航模型在边缘设备上的部署优化。近年来，空间感知机器人系统（Space-Aware Robot System for Terrain Crossing via Vision-Language Model, SARO）\cite{saro2025}将视觉语言模型用于空间感知四足机器人系统，增强了复杂地形语义理解和通过性判断能力；其重点更偏向语义地形理解，而不是以目标图像为条件的实时局部轨迹生成。

综上，国内外研究已经在视觉 SLAM、视觉目标导航、动作扩散策略和四足运动控制方面取得了大量进展，但面向“室内狭窄、弱先验环境中的四足机器人视觉目标导航”仍存在进一步研究空间。一方面，通用视觉导航模型多在轮式平台或离线数据集上验证，真实四足平台带来的相机抖动、运动接口和安全闭环问题并未被充分讨论；另一方面，四足机器人研究多集中在底层运动控制或复杂地形通过问题上，对扩散式视觉目标导航模型如何在四足平台上实现实时、安全、可复现的闭环部署关注不足。本文即围绕这一交叉问题展开研究。

\section{本文主要研究内容}

针对四足机器人在室内狭窄和弱先验环境中的视觉目标导航问题，本文围绕“算法推理优化、分层闭环系统、仿真与真机验证”三条主线开展研究，主要内容如下：

\begin{enumerate}
    \item 面向资源受限部署的统一推理模块构建。本文以 NoMaD 为基础建立视觉目标导航推理链路，围绕 DDIM 采样步数、Classifier-Free Guidance 条件引导、Test-Time Scaling 候选筛选和视觉编码器替换开展对比实验，分析不同优化方式对推理速度、轨迹质量和运行稳定性的影响，为扩散式导航模型在 Jetson Orin 等边缘设备上的实时运行提供依据。
    \item 面向 Lite3 四足机器人的分层闭环系统设计。本文构建统一推理模块、状态机、导航主机、中层控制映射和 Lite3 平台接口组成的系统架构，将高层生成的航点转换为可执行速度指令，并通过 UDP 通信下发至运动主机。同时，系统引入运行状态管理、异常输出抑制和安全停机逻辑，使视觉导航策略能够与真实四足平台形成稳定闭环。
    \item 面向仿真与真机的完整链路验证。本文在 MuJoCo 环境中搭建四足视觉导航验证流程，并在 Jetson Orin 与 Lite3 真机平台上完成相机输入、拓扑地图子目标选择、局部航点生成、速度指令下发和安全停机测试。通过离线实验、仿真闭环和真实部署三个层次，验证所构建系统的可运行性和工程可复现性。
\end{enumerate}

\section{论文章节安排}

本文整体内容按照“问题提出—统一推理模块构建—系统闭环验证—真机部署—总结展望”的逻辑展开，共分为五章。各章节之间的逻辑结构关系如图~\ref{fig:chapter-logic} 所示。第一章首先明确研究场景与问题边界，第二章围绕高层视觉导航策略构建基于动作扩散策略的统一视觉目标导航推理模块，第三章进一步将算法接入四足机器人闭环系统并完成仿真验证，第四章面向真实硬件平台完成部署与实验，第五章在前文结果基础上总结全文并讨论后续研究方向。

\begin{figure}[htbp]
    \centering
    \includegraphics[width=0.8\linewidth]{structure.png}
    \caption{本文各章节逻辑结构关系}
    \label{fig:chapter-logic}
\end{figure}

第一章为绪论，介绍四足机器人视觉目标导航的研究背景与意义，梳理传统导航、视觉目标导航、动作扩散策略和四足真实部署等方向的国内外研究现状，并明确本文的研究问题、主要内容和章节结构。

第二章围绕基于动作扩散策略的统一视觉目标导航推理模块展开，首先给出视觉目标导航、拓扑地图、航点、扩散模型和 NoMaD 机制等必要背景，随后介绍数据集、实验设备、基线推理验证、DDIM、CFG、TTS 与视觉编码器替换等推理模块实验。

第三章介绍基于分层架构与 MuJoCo 仿真的四足机器人视觉导航闭环系统设计与验证，说明轮式平台与足式平台的部署差异、Lite3 平台基础、统一推理模块、状态机、导航主机、中层控制映射、MuJoCo 仿真平台和闭环实验结果。

第四章给出 Jetson Orin 与 Lite3 真机部署流程，说明真机系统双链路、Orin 部署架构、真实拓扑地图构建、滑动窗口导航、速度指令下发、安全机制和真实环境实验结果。

第五章总结全文工作，归纳本文在算法优化、系统集成和真机验证方面的成果，并讨论后续可进一步拓展的研究方向。

\section{本章小结}

本章首先从具身智能和真实应用需求出发，说明了四足机器人视觉目标导航在消防巡检、室内导览和受限空间通行等场景中的研究价值，并进一步分析了视觉目标表达、四足平台形态和动作扩散策略带来的机遇与挑战。随后，本章按照导航建模、目标表达、扩散策略和四足部署的技术逻辑，对国内外相关研究进行了梳理，指出现有工作在真实四足平台上的扩散式视觉目标导航部署方面仍存在研究空间。最后，本章明确了本文面向 Lite3 四足机器人的研究内容和章节安排，为后续章节展开奠定基础。
